
\documentclass[12pt,a4paper]{ctexart}

\usepackage[a4paper,margin=1.8cm]{geometry}
\usepackage{array,booktabs,longtable,tabularx,multirow}
\usepackage{enumitem}
\usepackage{xcolor}
\usepackage{hyperref}
\usepackage{fancyhdr}
\usepackage{setspace}
\usepackage{titlesec}
\usepackage{lastpage}

\hypersetup{
    colorlinks=true,
    linkcolor=black,
    urlcolor=blue
}

\setlength{\parindent}{0pt}
\setlength{\parskip}{4pt}
\renewcommand{\arraystretch}{1.35}
\onehalfspacing

\pagestyle{fancy}
\fancyhf{}
\fancyhead[L]{RIS需求记录表}
\fancyhead[R]{DeliveraX / AI Agent需求流转}
\fancyfoot[C]{第 \thepage\ 页 / 共 \pageref{LastPage} 页}

\titleformat{\section}{\large\bfseries}{\thesection}{0.6em}{}
\titleformat{\subsection}{\normalsize\bfseries}{\thesubsection}{0.6em}{}

\newcolumntype{Y}{>{\raggedright\arraybackslash}X}
\newcolumntype{P}[1]{>{\raggedright\arraybackslash}p{#1}}

\begin{document}

\begin{center}
    {\LARGE \bfseries RIS（智能超表面）需求记录表模板}\\[0.5em]
    {\large Requirement Record Form for Reconfigurable Intelligent Surface}\\[0.8em]
    {\small 版本：V1.0 \quad 文档状态：草案 / 评审中 / 已基线（勾选其一）}
\end{center}

\vspace{0.6em}

%========================
% 文档基础信息
%========================
\section{文档基础信息}

\begin{tabularx}{\textwidth}{|P{3.2cm}|Y|P{3.2cm}|Y|}
\hline
文档编号 & % FILL_ME: 请Agent填写文档编号 (例如: RR-RIS-2026-001)
 & 项目名称 & % FILL_ME: 请Agent填写项目名称
 \\ \hline
所属客户 & % FILL_ME: 请Agent填写客户名称
 & 需求提出日期 & % FILL_ME: 请Agent填写日期 (格式: YYYY-MM-DD)
 \\ \hline
需求来源 & % FILL_ME: 请Agent填写需求来源
 & 当前阶段 & % FILL_ME: 请Agent填写当前项目阶段
 \\ \hline
责任人 & % FILL_ME: 请Agent填写责任人姓名
 & 参与人 & % FILL_ME: 请Agent填写参与人员列表
 \\ \hline
保密等级 & % FILL_ME: 请Agent填写保密等级
 & 优先级 & % FILL_ME: 请Agent填写优先级 (例如: P1)
 \\ \hline
关联文档 & % FILL_ME: 请Agent填写关联文档列表
 & 版本号 & % FILL_ME: 请Agent填写版本号 (例如: V1.0)
 \\ \hline
\end{tabularx}

\vspace{0.8em}

\textbf{填写说明：}
\begin{itemize}[leftmargin=2em]
    \item 本表用于客户需求采集、需求澄清、技术分析和立项评估的统一记录。
    \item 每条需求应尽量满足：唯一、明确、可验证、可追踪、可变更管理。
    \item 建议 AI Agent 从本表中自动抽取：场景、KPI、约束、接口、风险、验收项、责任人、待澄清项。
\end{itemize}

%========================
% 业务背景
%========================
\section{业务背景与建设目标}

\begin{tabularx}{\textwidth}{|P{3.2cm}|Y|}
\hline
业务场景 & % FILL_ME: 请Agent根据会议纪要填写具体的物理环境 (如: 地下二层停车场、长走廊等)
 \\ \hline
建设动因 & % FILL_ME: 请Agent总结现网存在的具体问题 (如: RSRP<-110dBm, 切换失败率高等)
 \\ \hline
建设目标 & % FILL_ME: 请Agent填写预期的技术指标 (如: RSRP>-105dBm, SINR>0dB等)
 \\ \hline
目标用户 & % FILL_ME: 请Agent推断受影响的人群 (如: 停车车主, 物流人员, 访客等)
 \\ \hline
业务价值 & % FILL_ME: 请Agent总结项目价值 (如: 降低投诉, 提升品牌形象, 节省CAPEX等)
 \\ \hline
\end{tabularx}

%========================
% 场景信息
%========================
\section{站点与场景信息}

\begin{tabularx}{\textwidth}{|P{3.2cm}|Y|P{3.2cm}|Y|}
\hline
建设地点 & % FILL_ME: 请Agent填写具体地理位置 (如: 某城市综合体地下停车场 B2-B3 层)
 & 场景类型 & % FILL_ME: 请Agent推断场景物理属性 (如: 室内 / 半封闭 / 多径效应显著)
 \\ \hline
覆盖面积 & % FILL_ME: 请Agent提取面积数值 (如: 约 6,000--10,000 平方米)
 & 层高 & % FILL_ME: 请Agent提取高度数值 (如: 约 3.6 米)
 \\ \hline
主要材质 & % FILL_ME: 请Agent总结环境材质 (如: 混凝土立柱、金属管道、高损耗墙体)
 & 目标区域 & % FILL_ME: 请Agent识别重点覆盖区 (如: 车道拐角、出入口、缴费区)
 \\ \hline
现网站点 & % FILL_ME: 请Agent填写现有覆盖设施 (如: 室分 / 宏站外打 / 无覆盖)
 & 传输条件 & % FILL_ME: 请Agent推断回传条件 (如: 机房可接入 / POE可用 / 光纤可达)
 \\ \hline
供电条件 & % FILL_ME: 请Agent推断供电方式 (如: AC 220V / 弱电井取电 / POE 供电)
 & 安装约束 & % FILL_ME: 请Agent提取现场限制 (如: 吊顶受限、消防要求、承重限制)
 \\ \hline
勘查状态 & % FILL_ME: 请Agent填写当前进度 (如: 已初勘 / 待复勘)
 & 附件 & % FILL_ME: 请Agent列出相关文件 (如: 平面图、照片、热力图)
 \\ \hline
\end{tabularx}

%========================
% 需求干系人
%========================
\section{干系人人与职责}

\begin{tabularx}{\textwidth}{|P{3.2cm}|P{3.0cm}|Y|P{3.0cm}|}
\hline
角色 & 姓名/部门 & 职责 & 联系方式 \\
\hline
客户负责人 & % FILL_ME: 请Agent提取客户方负责人姓名及部门
 & 需求确认、资源协调、验收签字 & % FILL_ME: 请Agent提取客户方联系电话
 \\ \hline
客户经理 & % FILL_ME: 请Agent提取我方客户经理姓名
 & 需求录入、客户沟通、版本流转 & % FILL_ME: 请Agent提取我方经理电话
 \\ \hline
产品经理 & % FILL_ME: 请Agent提取产品经理姓名
 & 需求分析、优先级判定、方案边界确认 & % FILL_ME: 请Agent提取产品经理电话
 \\ \hline
技术经理 & % FILL_ME: 请Agent提取技术负责人姓名
 & 技术可行性评估、参数建议、KPI拆解 & % FILL_ME: 请Agent提取技术负责人电话
 \\ \hline
交付经理 & % FILL_ME: 请Agent提取交付经理姓名
 & 施工计划、资源排期、验收组织 & % FILL_ME: 请Agent提取交付经理电话
 \\ \hline
\end{tabularx}

%========================
% 需求总览
%========================
\section{需求总览}

\begin{tabularx}{\textwidth}{|P{3.2cm}|Y|}
\hline
需求标题 & % FILL_ME: 请Agent提炼核心需求主题 (如: 地下停车场弱覆盖区域引入 RIS 进行定向覆盖增强)
 \\ \hline
需求摘要 & % FILL_ME: 请Agent总结客户核心诉求与预期效果 (如: 改善拐角/深处覆盖, 提升业务连续性)
 \\ \hline
需求分类 & % FILL_ME: 请Agent根据上下文判断类型 (如: 新建需求 / 网络优化 / 场景创新试点)
 \\ \hline
需求等级 & % FILL_ME: 请Agent评估项目紧迫性与复杂度 (如: 商务紧急度高; 技术复杂度中高)
 \\ \hline
期望上线时间 & % FILL_ME: 请Agent提取目标日期 (如: 2026-06-30)
 \\ \hline
验收方式 & % FILL_ME: 请Agent提取验收标准 (如: 现场测试 + 报表核验 + 客户签字确认)
 \\ \hline
\end{tabularx}

%========================
% 变更记录
%========================
\section{变更记录}

\begin{longtable}{|>{\centering\arraybackslash}p{1.4cm}|
                  >{\centering\arraybackslash}p{2.1cm}|
                  >{\centering\arraybackslash}p{1.8cm}|
                  >{\raggedright\arraybackslash}p{6.0cm}|
                  >{\centering\arraybackslash}p{1.8cm}|}
\hline
\textbf{版本} & \textbf{日期} & \textbf{修改人} & \textbf{修改说明} & \textbf{审批状态} \\
\hline
\endfirsthead

\hline
\textbf{版本} & \textbf{日期} & \textbf{修改人} & \textbf{修改说明} & \textbf{审批状态} \\
\hline
\endhead

\hline
\endfoot

\hline
\endlastfoot

% FILL_ME: 请Agent根据JSON中的ChangeLog列表填充以下行
% 格式示例: V0.1 & 2026-04-16 & 张三 & 初版创建，完成客户需求录入。 & 草案 \\
% 如果没有变更记录，可保留此行或删除
% 初始记录建议:
V1.0 & % FILL_ME: 请Agent填写当前日期 & % FILL_ME: 请Agent填写生成者(如: Agent/姓名) & % FILL_ME: 请Agent填写生成说明 (如: 初始需求生成) & % FILL_ME: 请Agent填写状态 (如: 草案) \\ \hline
\end{longtable}

%========================
% 审批签字
%========================
\section{审批与签字}

\begin{tabularx}{\textwidth}{|P{3.2cm}|Y|P{3.2cm}|Y|}
\hline
客户负责人签字 & \vspace{1.2cm} & 日期 & \\
\hline
客户经理签字 & \vspace{1.2cm} & 日期 & \\
\hline
产品经理签字 & \vspace{1.2cm} & 日期 & \\
\hline
技术经理签字 & \vspace{1.2cm} & 日期 & \\
\hline
交付经理签字 & \vspace{1.2cm} & 日期 & \\
\hline
\end{tabularx}

\vspace{1em}
\textbf{备注：}
\begin{itemize}[leftmargin=2em]
    \item 本模板中的指标与示例内容为演示填充，正式项目应以客户现网、勘查结果、方案评审结论和验收口径为准。
    \item 建议将本表作为“需求唯一输入源”，后续由 AI Agent 自动派生《需求分析报告》《方案评审单》《测试验收清单》《变更记录》。
\end{itemize}

\end{document}